\documentclass[11pt,a4paper]{article}

\usepackage[margin=1in]{geometry}
\usepackage{booktabs}
\usepackage{amssymb}
\usepackage{listings}
\usepackage{xcolor}
\usepackage{hyperref}
\usepackage{enumitem}
\usepackage{microtype}
\usepackage{array}

\hypersetup{
    colorlinks=true,
    linkcolor=blue!60!black,
    urlcolor=blue!60!black,
}

\lstdefinelanguage{Jimple}{
    keywords={staticinvoke,virtualinvoke,interfaceinvoke,specialinvoke,
              return,goto,if,new,newarray,instanceof,lengthof,
              static,public,private,protected,void,int,boolean,
              class,extends,implements},
    sensitive=true,
    comment=[l]{//},
    string=[b]",
    morekeywords={},
}

\lstdefinestyle{code}{
    basicstyle=\ttfamily\small,
    keywordstyle=\color{blue!70!black}\bfseries,
    commentstyle=\color{gray!70!black}\itshape,
    stringstyle=\color{red!70!black},
    frame=single,
    framesep=5pt,
    rulecolor=\color{gray!40},
    breaklines=true,
    showstringspaces=false,
    xleftmargin=10pt,
    xrightmargin=5pt,
    aboveskip=6pt,
    belowskip=6pt,
}

\lstset{style=code}

\title{\vspace{-2em}Monomorphization Optimizer --- Report}
\author{Kamyak Channa}
\date{April 2026}

\begin{document}
\maketitle
\vspace{-1.5em}

\section{Optimization Description}

\subsection*{What was implemented}

The optimizer is a Soot-based whole-program transformer that
\emph{devirtualizes} Java virtual and interface method calls when an
interprocedural type-flow analysis can prove that the receiver at a given
call site is always of exactly one concrete class. A proven-monomorphic
call is rewritten from \texttt{invokevirtual}/\texttt{invokeinterface} to
\texttt{invokestatic}, dispatched to a freshly synthesized static clone of
the concrete target method.

The baseline transformer already handled calls where the receiver flowed
through locals, parameters, casts, \texttt{new}, and invoke-results. This
work extends both the analysis and the rewriter along three dimensions
that were previously left conservative:

\begin{enumerate}[topsep=2pt,itemsep=0pt]
  \item \textbf{Instance and static field references} --- types written
    through a \texttt{FieldRef} are recorded into a global per-field map,
    and reads from fields populate the local's type set from that map.
  \item \textbf{Array references} --- writes to \texttt{ArrayRef}
    contribute to a global map keyed by the \emph{element type}, and
    reads produce that set.
  \item \textbf{Interface invokes} --- both \texttt{VirtualInvokeExpr} and
    \texttt{InterfaceInvokeExpr} are now analyzed and rewritten.
\end{enumerate}

\subsection*{Why it is useful}

Virtual dispatch in Java is expensive for three reasons: a hidden
vtable/itable lookup on every call, a cache-unfriendly memory
indirection, and the fact that the JIT cannot inline an unresolved
virtual callee without a guard. Static analysis that collapses a call to
a single concrete target turns the call into a plain \texttt{invokestatic}
--- faster to dispatch even in the pure interpreter, and structurally
inlineable once the JIT sees it. In practice, many Java programs use
\texttt{List<T>}, strategy objects held in fields, or arrays of a single
concrete subclass; without field/array tracking, the analysis is blind to
those shapes, so the original optimizer missed large classes of common
patterns (e.g.\ fluent \texttt{StringBuilder.append} chains devirtualize
only once field types are known).

\section{Analyses Used}

\subsection*{Type-flow analysis}

The core abstraction is a forward dataflow analysis over Soot's Jimple
CFG:

\begin{itemize}[topsep=2pt,itemsep=0pt]
  \item \textbf{Lattice element (\texttt{TypeState}):}
    \texttt{Map<Local, Set<SootClass>>}. An absent key means ``no type
    information'' (conservative $\bot$); an empty set means the local is
    known to be untyped. Join is point-wise set union.
  \item \textbf{Transfer function} (\texttt{TypeFlowAnalysis.flowThrough}):
    \texttt{IdentityStmt} seeds parameters from the calling context;
    \texttt{NewExpr} produces a singleton; \texttt{Local := Local}
    propagates; \texttt{CastExpr} narrows via
    \texttt{FastHierarchy.canStoreClass}; \texttt{InvokeExpr} consumes the
    callee's return-type summary for the current arg context;
    \texttt{FieldRef} and \texttt{ArrayRef} read from the global maps;
    writes to those locations update the globals.
\end{itemize}

\subsection*{Dimensions of precision}

\begin{itemize}[topsep=2pt,itemsep=0pt]
  \item \textbf{Flow-sensitive.} Type sets are maintained per program
    point, computed by Soot's \texttt{ForwardFlowAnalysis}.
  \item \textbf{Context-sensitive on method parameters} (1-CFA-like). The
    analyzer is parameterised by a \texttt{paramCtx: Map<Integer,
    Set<SootClass>>} keyed by parameter index ($-1$ for the receiver).
    The worklist enqueues \texttt{(method, ctx)} pairs; method summaries
    are $\mathit{ctx} \to \mathit{returnTypes}$. Virtual targets are
    enqueued with the context refined to the specific receiver type,
    which is how per-dispatch-target precision is recovered.
  \item \textbf{Field-based} (not field-sensitive). A field is treated as
    a single abstract cell shared across all instances ---
    \texttt{SootField} $\to$ \texttt{Set<SootClass>}. This trades some
    precision for simplicity and scales well to large heaps.
  \item \textbf{Array-element-type-based.}
    \texttt{Type} $\to$ \texttt{Set<SootClass>}, keyed by the declared
    element type (\texttt{ArrayRef.getType()}). All arrays with the same
    element type are merged.
  \item \textbf{Conservative on unmodelled rhs.} Any assignment whose
    right-hand side the analysis does not understand (constants,
    intrinsics it doesn't model) \emph{clears} the lhs set, so stale
    types from a previous store can never survive.
\end{itemize}

\subsection*{Assumptions}

\begin{itemize}[topsep=2pt,itemsep=0pt]
  \item SPARK call graph, built on-the-fly
    (\texttt{-p cg.spark on-fly-cg:true}).
  \item Reflection is resolved via the TamiFlex \texttt{refl.log} that
    accompanies each DaCapo input directory.
  \item Only application classes (as seen by Soot, which --- with
    \texttt{-app} --- includes TamiFlex-dumped JDK classes) are analyzed.
  \item \textbf{JDK rewriting is disabled.} \texttt{\$mono} clones are
    only created when the concrete target's declaring class is outside
    \texttt{java.*}, \texttt{javax.*}, \texttt{sun.*}, \texttt{jdk.*},
    \texttt{com.sun.*}. See \S3 for why.
  \item Correctness hinges on \textbf{outer fixed-point iteration}: the
    field/array global maps are re-checked after each worklist pass, and
    a new pass is started if either map grew. Monotonicity (maps only
    grow) guarantees termination in a small number of passes.
\end{itemize}

\section{Transformation Details}

\subsection*{Before / after snippets}

\textbf{Field-flowing call (interface dispatch through an instance field).}
Before monomorphization (excerpt of \texttt{Kennel.call}):

\begin{lstlisting}[language=Jimple]
$r1 = r0.<Kennel: Animal resident>;
$r2 = interfaceinvoke $r1.<Animal: String speak()>();
return $r2;
\end{lstlisting}

After:

\begin{lstlisting}[language=Jimple]
$r1 = r0.<Kennel: Animal resident>;
$mono_cast_0 = (Dog) $r1;
$r2 = staticinvoke <Dog: String speak$mono$Dog(Dog)>($mono_cast_0);
return $r2;
\end{lstlisting}

The \texttt{interfaceinvoke} is replaced by an \texttt{invokestatic} to a
synthesized \texttt{speak\$mono\$Dog}, a static copy of \texttt{Dog.speak}
whose first parameter is the receiver. An explicit checkcast is inserted
because \texttt{\$r1}'s declared type (\texttt{Animal}) is wider than the
static clone's first parameter type (\texttt{Dog}).

\textbf{Array-element-flowing call.} Before:

\begin{lstlisting}[language=Jimple]
r5 = r0[i3];             // r0: Shape[]
$i1 = interfaceinvoke r5.<Shape: int area()>();
\end{lstlisting}

After:

\begin{lstlisting}[language=Jimple]
r5 = r0[i3];
$mono_cast_0 = (Square) r5;
$i1 = staticinvoke <Square: int area$mono$Square(Square)>($mono_cast_0);
\end{lstlisting}

\textbf{Generated static clone} (added as a new method on the concrete
class):

\begin{lstlisting}[language=Jimple]
static String speak$mono$Dog(Dog) {
    Dog r0;
    r0 := @parameter0: Dog;   // originally: @this: Dog
    return "woof";
}
\end{lstlisting}

Identity statements are fixed up when cloning: \texttt{@this: T} becomes
\texttt{@parameter0: T}, and each existing \texttt{@parameterN: U} is
shifted to \texttt{@parameter(N+1): U}.

\subsection*{Key challenges}

\begin{enumerate}[topsep=2pt,itemsep=2pt]
  \item \textbf{Cross-method field/array data flow.} Writers and readers
    of a field can be in different methods; the worklist visits methods
    in arbitrary order, so a reader may be analyzed before any writer
    has been seen. We solved this with a monotone outer fixed-point: the
    \texttt{(method, ctx)} worklist is re-seeded until the global field
    and array maps stop growing. In practice two or three passes suffice
    even on DaCapo-scale input, because the maps are small.

  \item \textbf{Field-sensitive vs.\ field-based.} True
    instance-field sensitivity requires pointer-analysis-grade heap
    abstractions. We deliberately chose the field-based abstraction (one
    abstract cell per field symbol). This loses precision when a field
    can legitimately hold different types in different instances, but it
    keeps the analysis fast and is sound --- any over-approximation
    leaves the call site conservatively virtual.

  \item \textbf{Interface invokes were silently ignored by the original
    rewriter.} Fixing fields/arrays without fixing interface invokes
    would have hidden the benefit, because typical field-loaded
    receivers are declared with an interface type (\texttt{Animal},
    \texttt{Op}, \texttt{Shape}). Both the analyser's receiver-type
    collection and the rewriter's instance-invoke dispatch were
    generalized to handle \texttt{InterfaceInvokeExpr}.

  \item \textbf{Receiver type mismatch after devirtualization.} A call
    site typed \texttt{Animal a; a.speak()} passes a local declared as
    \texttt{Animal}, but the clone's first parameter is declared as the
    concrete class (\texttt{Dog}). The JVM verifier rejects this
    mismatch. The rewriter now inserts a \texttt{CastExpr} assign into a
    fresh local immediately before the call when the base local's
    declared type is not equal to the concrete type.

  \item \textbf{JDK class rewriting breaks at runtime.} With \texttt{-app}
    Soot treats TamiFlex-dumped JDK classes as application classes, and
    the original rewriter happily synthesized \texttt{\$mono} clones on
    \texttt{java.lang.StringBuilder} and friends. Emitted classes live
    under \texttt{sootOutput/}, but at runtime the JVM loads
    \texttt{java.*} from the boot class-loader, not from our output
    directory --- so every rewritten call site would fail with
    \texttt{NoSuchMethodError}. The rewriter now skips any concrete
    target whose declaring class is under a JDK package prefix. The
    \emph{analysis} still traverses JDK methods, so JDK return types
    still refine application call sites; only the rewrite is suppressed.
\end{enumerate}

\section{Result Tables}

\subsection{Functional correctness}

Four purpose-built micro testcases exercise the new capabilities. For
each, the original bytecode (JVM baseline) and the monomorphized
bytecode (after applying the transformer, then running the emitted
\texttt{.class} files) are executed and their stdout compared.

\begin{table}[h]
\centering
\small
\begin{tabular}{@{}lllccc@{}}
\toprule
Test case            & Exercised feature                         & JVM   & After mono & Match & Sites \\
\midrule
\texttt{field\_case}        & interface dispatch via instance field      & woof  & woof       & \checkmark & 3 \\
\texttt{array\_case}        & interface dispatch via array element       & 29    & 29         & \checkmark & 1 \\
\texttt{mixed\_case}        & array held inside an instance field        & RR    & RR         & \checkmark & 5 \\
\texttt{static\_field\_case}& static field written/read across methods   & hello & hello      & \checkmark & 1 \\
\bottomrule
\end{tabular}
\caption{Functional correctness: all outputs match the JVM baseline.}
\end{table}

Disassembly (\texttt{javap -c}) confirms each devirtualized call site is
an \texttt{invokestatic} to a \texttt{\$mono} clone and not an
\texttt{invokeinterface}/\texttt{invokevirtual}.

\subsection{Performance}

\subsubsection*{Microbenchmarks (\texttt{-Xint}, $4\times 10^8$ iterations, median of 5)}

Unit: seconds wall-clock. Workload = tight loop dispatching through the
feature under test.

\begin{table}[h]
\centering
\small
\begin{tabular}{@{}lrrrrr@{}}
\toprule
Microbench         & Sites & Baseline (s) & Optimized (s) & $\Delta$ & Speedup \\
\midrule
\texttt{bench\_array} & 1  & 7.86  & 6.89  & $+12.3\%$ & $1.14\times$ \\
\texttt{bench\_mixed} & 2  & 8.50  & 8.56  & $-0.7\%$  & $0.99\times$ \\
\texttt{bench\_field} & 3  & 12.28 & 13.07 & $-6.4\%$  & $0.94\times$ \\
\bottomrule
\end{tabular}
\caption{Micro-benchmarks, pure interpreter mode.}
\end{table}

Under \texttt{-Xint} the only saving is
\texttt{invokeinterface}$\to$\texttt{invokestatic} (one fewer hash-table
lookup per call); the inserted \texttt{checkcast} costs about as much in
the pure interpreter. \texttt{bench\_array} is a clean win because the
dispatch sits in the hottest possible loop and the cast is cheap compared
to an interface lookup; \texttt{bench\_field} shows a small regression
because three rewritten sites include \texttt{Holder.run} whose original
\texttt{virtualinvoke} was already fast. The transformation is structured
to feed the JIT, which inlines the static clone and folds the cast away
--- that benefit is invisible under \texttt{-Xint}.

\subsubsection*{DaCapo 9.12 benchmarks (\texttt{-s default}, \texttt{-Xint}, 3 runs each)}

Unit: seconds wall-clock. ``Sites'' is the number of application-level
call sites the transformer rewrote (JDK targets deliberately skipped).

\begin{table}[h]
\centering
\small
\begin{tabular}{@{}lrrrrrr@{}}
\toprule
Benchmark & Sites & Baseline mean & Baseline min & Optimized mean & Optimized min & Speedup \\
\midrule
\texttt{xalan}   & 282   & 4.65\,s & 4.55\,s & 4.19\,s & 4.15\,s & $\mathbf{1.11\times}$ \\
\texttt{avrora}  & 1\,090 & 8.70\,s & 8.32\,s & 8.40\,s & 8.32\,s & $1.04\times$ \\
\texttt{fop}     & 2\,231 & 3.22\,s & 3.19\,s & 3.18\,s & 3.13\,s & $1.01\times$ \\
\texttt{luindex} & 1\,279 & 5.16\,s & 5.05\,s & 5.19\,s & 5.06\,s & $0.99\times$ \\
\bottomrule
\end{tabular}
\caption{DaCapo 9.12 benchmarks, JIT disabled. \texttt{batik} excluded
because DaCapo-9.12 still requires \texttt{com.sun.image.codec.jpeg},
which was removed in JDK 9.}
\end{table}

(Per-run times, logs, and transformed class files are archived under
\texttt{bench-results/}.)

\subsection*{Discussion}

\begin{itemize}[topsep=2pt,itemsep=0pt]
  \item \texttt{xalan} sees the largest genuine win ($+11\%$) despite
    having the fewest rewritten sites: those 282 sites sit on the tight
    XSLT template-dispatch loop that dominates its interpreted runtime.
  \item \texttt{luindex}'s interpreted time is spent mostly in JDK code
    paths (file channels, regular expressions,
    \texttt{StringBuilder}) --- exactly the surface our JDK-skip filter
    refuses to rewrite, so there is nothing to gain.
  \item \texttt{avrora} and \texttt{fop} show small positive deltas
    ($+1$--$4\%$); their hot paths are mixed app/JDK, so devirtualization
    dilutes proportionally.
  \item No benchmark regresses beyond $\sim1\%$ noise, indicating the
    checkcast-insertion cost and the extra static-frame overhead are
    well-contained.
\end{itemize}

\subsection{Reproduction}

All numbers can be reproduced with the supplied driver script:

\begin{lstlisting}[language=bash]
# default: luindex avrora xalan fop, 3 runs, -s default, JIT off
./bench.sh

# quick smoke test
SIZE=small RUNS=1 ./bench.sh

# subset with more runs
BENCHES="xalan avrora" RUNS=5 ./bench.sh
\end{lstlisting}

The script recompiles the transformer, applies it once per benchmark to
emit a \texttt{mono-<bench>/} directory of \texttt{.class} files, then
runs \texttt{java -Xint -cp dacapo.jar Harness -s default <bench>}
\texttt{RUNS} times for the baseline and again with
\texttt{mono-<bench>/} prepended to the classpath, emitting per-run
times, means, mins, and speedup ratios to stdout and
\texttt{bench-results/}.

\end{document}
